\begin{thema}{Β}
Δίνεται η γραφική παράσταση $C_f$ μιας συνεχούς $f:[-2, 3]\to\mathbb{R}$, παραγωγίσιμης στο $(-2,2)$, μαζί μe τις εφαπτόμενές της σε δύο σημεία $A(0, 2)$ και $B(2, 0)$. Η εφαπτομένη στο $A$ έχει κλίση $1$ και στο $B$ κλίση $-2$.

\begin{center}
\begin{tikzpicture}[scale=0.8]
  \draw[help lines, gray!30] (-3,-2) grid (4,4);
  \draw[-latex] (-3.3,0) -- (4.3,0) node[right] {$x$};
  \draw[-latex] (0,-2.3) -- (0,4.3) node[above] {$y$};
  \foreach \x in {-2,-1,1,2,3}
    \draw (\x,0.1) -- (\x,-0.1) node[below,font=\footnotesize] {$\x$};
  \foreach \y in {-1,1,2,3}
    \draw (0.1,\y) -- (-0.1,\y) node[left,font=\footnotesize] {$\y$};
  % curve
  \draw[very thick, maincolor!70!black, smooth, tension=0.5]
    plot coordinates {(-2,0) (-1,0.8) (0,2) (1,2.5) (2,0) (3,-1)};
  % tangent at A(0,2), slope 1
  \draw[dashed, secondcolor] (-1,1) -- (1.5,3.5);
  \node[secondcolor,font=\footnotesize] at (1.5,3.8) {$\varepsilon_A$};
  % tangent at B(2,0), slope -2
  \draw[dashed, secondcolor] (1,2) -- (3,-2);
  \node[secondcolor,font=\footnotesize] at (3,-1.7) {$\varepsilon_B$};
  \filldraw[maincolor!70!black] (0,2) circle (2pt) node[above left,font=\footnotesize] {$A(0,2)$};
  \filldraw[maincolor!70!black] (2,0) circle (2pt) node[below left,font=\footnotesize] {$B(2,0)$};
\end{tikzpicture}
\end{center}

\begin{erwthma}
  \item Να εφαρμόσετε το Θεώρημα \eng{Rolle} στο $[-2, 2]$ και να αποδείξετε ότι υπάρχει $\xi \in (-2, 2)$ με $f'(\xi) = 0$. Ποια η γεωμετρική ερμηνεία του $\xi$ στο σχήμα?
  \item Γράψτε τις εξισώσεις των εφαπτομένων $\varepsilon_A$ και $\varepsilon_B$.
  \item Αξιοποιώντας τα δεδομένα του σχήματος, να υπολογίσετε τα όρια:
  \[ \lim\limits_{x \to 0} \frac{f(x) - 2}{x} \qquad \text{και} \qquad \lim\limits_{x \to 2} \frac{f(x)}{x - 2}. \]
  \item Αν $h(x) = f(x) \cdot x$, να βρείτε $h'(0)$ και $h'(2)$.
\end{erwthma}
\end{thema}
